\documentclass[12pt,a4paper]{article}

% ------------------------------------------------
% Packages
% ------------------------------------------------
\usepackage[utf8]{inputenc}
\usepackage[T1]{fontenc}
\usepackage[german]{babel}

\usepackage{geometry}
\geometry{margin=2.5cm}

\usepackage{longtable}
\usepackage{booktabs}
\usepackage{array}
\usepackage{graphicx}
\usepackage[hidelinks]{hyperref}
\usepackage{parskip}
\usepackage{xcolor}

% ------------------------------------------------
% Header / Footer
% ------------------------------------------------
\usepackage{fancyhdr}
\pagestyle{fancy}
\fancyhf{}
\lhead{Projektplan}
\rhead{\thepage}
\geometry{a4paper, margin=1in}

% ------------------------------------------------
% Document
% ------------------------------------------------
\begin{document}

% ------------------------------------------------
% Deckblatt
% ------------------------------------------------
\begin{titlepage}
    \centering

    \vspace*{3cm}

    {\Huge \textbf{PROJEKTPLAN}}\\[1.5cm]

    {\LARGE Workflowbasierte Rückbestätigung}\\[0.3cm]
    {\LARGE von Lieferterminen}\\[3cm]

    {\Large Projektarbeit}\\[2cm]

    {\large
    MINOVA / THWS
    }\\[1cm]

    {\large
    \today
    }

    \vfill

\end{titlepage}

% ------------------------------------------------
\tableofcontents
\newpage

% ------------------------------------------------
\section{Einleitung}

Im aktuellen DISPO-Prozess werden Liefertermine bzw. Lieferoptionen manuell mit Kunden abgestimmt. 
Dieser Ablauf ist zeitintensiv, fehleranfällig und bei steigender Auftragszahl schwer skalierbar.

Ziel des Projekts ist die Entwicklung eines prototypischen Systems zur workflowbasierten Rückbestätigung von Lieferterminen. 
Das System unterstützt die Disposition bei der strukturierten Einholung und Verarbeitung von Kundenrückmeldungen zu vorgeschlagenen Lieferoptionen.

Die Prozesslogik wird mithilfe von BPMN 2.0 modelliert und durch die Workflow-Engine Camunda ausgeführt.

% ------------------------------------------------
\section{Projektziel}

Das Projekt verfolgt folgende Ziele:

\begin{itemize}
    \item Erfassung und Anzeige relevanter Bestell- und Auftragsdaten
    \item Darstellung von durch die Disposition vorgeschlagenen Lieferoptionen
    \item Digitale Kommunikation mit Kunden
    \item Verarbeitung von Kundenreaktionen wie Bestätigung, Ablehnung oder ausbleibende Antwort
    \item Workflowgesteuerte Prozessausführung mit Camunda
    \item Übergabe bestätigter Aufträge an die Disposition
    \item Verbesserung der Transparenz und Nachverfolgbarkeit
\end{itemize}

% ------------------------------------------------
\section{Projektphasen und Meilensteine}

\begin{longtable}{p{4.5cm} p{6cm} p{3cm} p{2cm}}
\toprule
\textbf{Phase} & \textbf{Beschreibung} & \textbf{Zeitraum} & \textbf{Meilenstein} \\
\midrule

P1: Projektdefinition und Analyse 
& Projektidee abstimmen, Anforderungen analysieren und Projektumfang definieren 
& 28.04--04.05 
& M1 \\

P2: Planung und Grundstruktur des Systems 
& Lastenheft, Pflichtenheft und Projektplan erstellen sowie erste Backend- und Frontendstruktur aufsetzen 
& 05.05--11.05 
& M2 \\

P3: Architektur und Schnittstellen 
& Systemarchitektur, DISPO-Schnittstelle, Datenmodell und Backend-Grundstruktur definieren 
& 12.05--25.05 
& M3 \\

P4: Weiterentwicklung von Backend und Workflow 
& Backend, Frontend und Camunda-Workflow funktional erweitern 
& 26.05--22.06 
& M4 \\

P5: End-to-End-Prozess 
& Durchgängigen Prozess von Auftrag über Kundenrückmeldung bis Statusänderung umsetzen 
& 23.06--06.07 
& M5 \\

P6: Testphase und Fehlerbehebung 
& Systemtests durchführen, Fehler identifizieren und beheben 
& 07.07--20.07 
& M6 \\

P7: Finalisierung des Prototyps 
& Prototyp stabilisieren, Demo vorbereiten und finale Anpassungen durchführen 
& 21.07--03.08 
& M7 \\

P8: Dokumentation und Präsentation 
& Projektdokumentation und Abschlusspräsentation erstellen 
& 04.08--16.09 
& M8 \\

\bottomrule
\end{longtable}

% ------------------------------------------------
\section{Meilensteine}

\begin{longtable}{p{3cm} p{11cm}}
\toprule
\textbf{Meilenstein} & \textbf{Beschreibung} \\
\midrule

M1 & Projekt bestätigt und Projektumfang definiert \\

M2 & Lastenheft, Pflichtenheft und Projektplan abgeschlossen; erste Version von Backend und Frontend erstellt \\

M3 & Systemarchitektur und DISPO-Schnittstelle definiert; Backend-Grundstruktur implementiert \\

M4 & Backend, Frontend und Workflow erweitert \\

M5 & End-to-End-Prozess funktionsfähig \\

M6 & Tests abgeschlossen und Fehler behoben \\

M7 & Prototyp abgeschlossen \\

M8 & Dokumentation und Abschlusspräsentation abgeschlossen \\

\bottomrule
\end{longtable}

% ------------------------------------------------
\section{Zeitplan}

Der Projektzeitplan orientiert sich an den definierten Projektphasen und Meilensteinen. 
Die einzelnen Phasen bauen aufeinander auf und führen schrittweise zur Fertigstellung des Prototyps sowie zur Abschlusspräsentation.

\begin{figure}[h!]
    \centering
    \includegraphics[width=\textwidth]{assets/projectplan.png}
    \caption{Projektphasen und Meilensteine}
    \label{fig:projektplan}
\end{figure}

% ------------------------------------------------
\section{Systemarchitektur}

Das System besteht aus folgenden Komponenten:

\begin{itemize}
    \item Frontend (Web UI)
    \item Backend-Service
    \item Camunda Workflow Engine
    \item PostgreSQL Datenbank
    \item E-Mail Kommunikationsschnittstelle
\end{itemize}

% ------------------------------------------------
\section{Technologien}

\begin{itemize}
    \item BPMN 2.0
    \item Camunda
    \item REST APIs
    \item PostgreSQL
    \item Java oder Node.js
    \item HTML/CSS/JavaScript
\end{itemize}

% ------------------------------------------------
\section{Rollen}

\begin{itemize}
    \item Disponent
    \item Kunde
    \item System
\end{itemize}

% ------------------------------------------------
\section{Qualitätssicherung}

Zur Sicherstellung der Qualität werden folgende Maßnahmen durchgeführt:

\begin{itemize}
    \item Funktionstests
    \item Integrationstests
    \item Prozesssimulationen
    \item Validierung der BPMN-Prozesse
    \item Fehleranalyse und Fehlerbehebung
\end{itemize}

% ------------------------------------------------
\section{Akzeptanzkriterien}

Das Projekt gilt als erfolgreich abgeschlossen, wenn:

\begin{itemize}
    \item Kundenreaktionen korrekt verarbeitet werden
    \item Statusänderungen nachvollziehbar gespeichert werden
    \item der End-to-End-Prozess funktionsfähig ist
    \item BPMN-Prozesse in Camunda erfolgreich ausgeführt werden
    \item Frontend und Backend korrekt kommunizieren
\end{itemize}

% ------------------------------------------------
\section{Erwartetes Ergebnis}

Am Ende des Projekts soll ein funktionsfähiger Prototyp eines workflowbasierten Systems zur Rückbestätigung von Lieferterminen vorliegen. 

Das System soll Kundenrückmeldungen strukturiert verarbeiten und bestätigte Aufträge transparent für die Disposition bereitstellen.

% ------------------------------------------------
\end{document}